\documentclass[10pt,a4paper]{amsart}
\usepackage[margin=19mm]{geometry}
\usepackage{amsmath,amssymb,mathtools}
\usepackage[scr=rsfs,cal=euler]{mathalfa}
\usepackage{xfrac}
\usepackage[unicode,hidelinks,bookmarks=false]{hyperref}
\newcommand{\ie}{i.e.\ }
\newcommand{\cf}{cf.\ }
\newcommand{\resp}{respectively\ }
\newcommand{\Subsection}{\subsection}
\providecommand{\nobreakdash}{\nobreak\hbox{-}}
\setlength{\emergencystretch}{2em}
\multlinegap0pt
\usepackage{bm}
\usepackage{bbm}
\usepackage{dsfont}
\usepackage{xspace}
\usepackage{cancel}
\usepackage{mathtools}
\usepackage{amsthm}
\makeatletter
\@ifundefined{lemma}{\newtheorem{lemma}{Lemma}}{}
\@ifundefined{proposition}{\newtheorem{proposition}{Proposition}}{}
\makeatother
\newcommand{\E}{\mathbb{E}}
\newcommand{\N}{\mathbb{N}}
\newcommand{\R}{\mathbb{R}}
\newcommand{\p}{\mathbb{P}}
\title[Character sums over smooth numbers]{Character sums over smooth numbers}
\author{Seth Hardy, Max Wenqiang Xu}
\date{Source: arXiv:2607.00592v1 (CC BY 4.0)}
\begin{document}
\maketitle

\noindent\emph{Source and licence.} Character sums over smooth numbers. Version arXiv:2607.00592v1 (CC BY 4.0), distributed under the Creative Commons Attribution 4.0 International licence, as recorded in the arXiv licence metadata for this exact version and shown on the abstract page https://arxiv.org/abs/2607.00592v1. Full rights notice: licenses/arxiv-2607.00592-CC-BY-4.0.txt.\par\smallskip

\noindent\textit{Editorial scope.} Counted unit: the Proposition whose source label is \texttt{p: L1 norm} in the article (source0.tex lines 226--231, source label \texttt{p: L1 norm}), delivered together with the article's own complete original proof of it (source0.tex lines 237--265, one \texttt{proof} environment of 29 lines). No mathematics is written, repaired or re-derived: statement and proof are the article's own text line for line. The proof is detached from the statement in the source: the article states the result at lines 226--231 and prints its proof at lines 237--265, and the proof environment's own header names the statement, so the association is the article's own. Required context retained uncounted: l: good generating function (lines 196--205); l: large value bound (lines 212--218); equ: u upper bound log zeta (lines 302--304). Every definition, display and label of those lines is retained unchanged. Omitted from this package and not redistributed: 1--195, 206--211, 219--225, 232--236, 266--301, 305--487. Nothing inside a retained excerpt is omitted or altered: every retained body is delivered exactly as the article's lines are written, so no omission receipt of any kind accompanies this package. Citations: the retained excerpts carry no citation command at all, so no bibliography entry of the article is transcribed and no reference is redistributed. Reference adaptation: none was needed and none was made; every reference inside the delivered unit resolves inside it, so no unresolved reference or citation is produced. Printed slips kept literal: no printed word, symbol, hypothesis or number of the counted statement or of its proof was altered. Layout and class: the article's own class is \texttt{amsart} with packages including geometry, inputenc, upgreek, amsmath, mathtools, esint, dsfont, amssymb, enumitem, biblatex, hyperref; the wrapper is the lane's pinned \texttt{amsart} editorial wrapper at 10pt on A4, which restates the theorem-like environments the retained text needs and adds the article's own macro definitions that the retained text needs (\texttt{\textbackslash E}, \texttt{\textbackslash N}, \texttt{\textbackslash R}, \texttt{\textbackslash p}), with extra wrapper packages (bm, bbm, dsfont, xspace, cancel, mathtools). The delivered numbering is the unit's own. Assets: no retained span contains a figure environment, a graphics-inclusion command, a PDF-page inclusion, a table or a TikZ picture; no image file is copied into this package, the package declares no asset dependency and no image file is redistributed with it. Licence: version arXiv:2607.00592v1 of this article is distributed under the Creative Commons Attribution 4.0 International licence, as recorded in the arXiv licence metadata for this exact version and shown on the abstract page https://arxiv.org/abs/2607.00592v1. A full rights notice is delivered at licenses/arxiv-2607.00592-CC-BY-4.0.txt. Characters: every retained excerpt is pure ASCII, so the delivered engine, XeLaTeX with its default Latin Modern text font, reports no missing character.
\begin{lemma}[Approximation by prime sum expansion]\label{l: good generating function}
Suppose that $(\log x)^4 \leq y \leq x$, let $u = \frac{\log x}{\log y}$, and let $\alpha = \alpha (x,y) $ be the saddle point. If $K \in \N$ is such that $x \leq y^K \leq \min \{ q, x^{3/2} \}$, then uniformly for $t \in \R$, we have
\[
\E_{\chi (q)} \biggl| \sum_{\substack{n \leq y^K \\ P(n) \leq y}} \frac{\chi(n)}{n^{\alpha / 2 + it}} - \sum_{k = 0}^{K} \frac{P (\alpha / 2 + it, \chi, y)^{k}}{k!} \biggr| \ll E (x, y, K), 
\]
where
\[
E (x, y, K) \coloneqq \zeta (\alpha, y)^{1/2} \biggl(y^{(1/2 - \alpha)/2} + \exp \biggl\{ - \frac{u}{8} \biggl( \frac{K}{u} -1 \biggr)^2 \biggr\} \biggr).
\]
\end{lemma}

\begin{lemma}[Tail bound for sum over primes]\label{l: large value bound}
Suppose that $(\log x)^4 \leq y \leq x^{\frac{1}{4 \log \log x}}$, let $u = \frac{\log x}{\log y}$, and let $\alpha = \alpha (x,y)$ be the saddle point. There exists an absolute constant $c \geq 1$ so that for $V \geq c \sqrt{u}$ where $u=\frac{\log x}{\log y}$, and for any $t \in \R$, we have
\[
\p_{\chi(q)} \bigl( |P(\alpha / 2 + it, \chi , y) | > V \bigr) \ll \sqrt{k} \biggl( \frac{k u}{2 V^2} \biggr)^{k} ,
\]
uniformly for $k \in \N$ satisfying $k \leq \lfloor \frac{\log q}{\log y} \rfloor$.
\end{lemma}

\begin{equation}\label{equ: u upper bound log zeta}
u \leq \log \zeta (\alpha, y) + c.
\end{equation}

\begin{proposition}\label{p: L1 norm}
Suppose that $x \leq q$ and $(\log x)^4 \leq y \leq x^{\frac{1}{4 \log \log x}}$. Let $\alpha = \alpha (x,y)$ be the saddle point, let $u = \frac{\log x}{\log y}$, and define $\mathcal{K} \coloneqq \lfloor \frac{\log \min \{ x^{3/2}, q \}}{\log y} \rfloor$. Uniformly for any $t \in \R$, we have
\[
\E_{\chi (q)} \biggl| \sum_{\substack{n \leq y^{\mathcal{K}} \\ P(n) \leq y}} \frac{\chi(n)}{n^{\alpha / 2 + it}} \biggr| \ll \zeta (\alpha, y)^{1/2} \biggl(y^{(1/2 - \alpha)/2} + \exp \biggl\{ - \frac{u}{8} \biggl( \frac{\log \min\{ q, x^{3/2} \}}{\log x} - 1 \biggr)^2 \biggr\} \biggr).
\]
\end{proposition}

\begin{proof}[Proof of Proposition~\ref{p: L1 norm}, assuming Lemmas~\ref{l: good generating function} and~\ref{l: large value bound}. ]
By linearity of expectation, the term that we wish to bound is equal to
\[
\E_{\chi (q)} \mathbf{1}_{|P (\alpha / 2 + it, \chi, y)| \leq V} \Bigl| \sum_{\substack{n \leq y^{\mathcal{K}} \\ P(n) \leq y}} \frac{\chi(n)}{n^{\alpha /2+it}} \Bigr| + \E_{\chi (q)} \mathbf{1}_{|P (\alpha / 2 + it, \chi, y)| > V} \Bigl| \sum_{\substack{n \leq y^{\mathcal{K}} \\ P(n) \leq y}} \frac{\chi(n)}{n^{\alpha /2+it}} \Bigr| .
\]
Applying the Cauchy--Schwarz inequality, the second term in the previous display is bounded above by
\[
\p_{\chi (q)} \bigl( |P (\alpha / 2 + it, \chi, y)| > V \bigr)^{1/2} \biggl( \E_{\chi (q)} \Bigl| \sum_{\substack{n \leq \min \{x^{3/2}, q\} \\ P(n) \leq y}} \frac{\chi(n)}{n^{\alpha / 2 + it}} \Bigr|^2 \biggr)^{1/2}
\]
By applying orthogonality and dropping the upper bound condition on the size of $n$, we see that the expectation inside the parentheses is $\ll \zeta(\alpha, y)$. To bound the probability, we apply Lemma~\ref{l: large value bound} with $V = u/3$ and $k = 2 u / 9e$, to find the previous display is
\[
\ll k^{1/4} \biggl( \frac{k u}{2 V^2} \biggr)^{k/2} \zeta(\alpha, y)^{1/2} \ll u^{1/4} e^{-u/9e} \zeta(\alpha, y)^{1/2} \ll e^{-u/25} \zeta(\alpha, y)^{1/2} .
\]
We note that there is considerable flexibility in these choices, and since the right-hand side of Proposition~\ref{p: L1 norm} is always $\gg \zeta (\alpha, y)^{1/2} e^{-u/32}$, this bound is acceptable. It remains to show that
\begin{multline*}
\E_{\chi (q)} \mathbf{1}_{|P (\alpha/2 + it, \chi, y)| \leq V} \Bigl| \sum_{\substack{n \leq y^{\mathcal{K}} \\ P(n) \leq y}} \frac{\chi(n)}{n^{\alpha / 2 + it}} \Bigr| \\ 
\ll \zeta (\alpha, y)^{1/2} \biggl(y^{(1/2 - \alpha)/2} + \exp \biggl\{ - \frac{u}{8} \biggl( \frac{\log \min\{ q, x^{3/2} \}}{\log x} -1 \biggr)^2 \biggr\} \biggr),
\end{multline*}
when $V = u/3$. By the triangle inequality, we can upper-bound the left-hand side by
\[
\E_{\chi(q)} \mathbf{1}_{|P (\alpha/2 + it, \chi, y)| \leq V} \Bigl| \sum_{k=0}^{\mathcal{K}} \frac{P (\alpha/2 + it, \chi, y)^k}{k!} \Bigr| + \E_{\chi (q)} \biggl| \sum_{\substack{n \leq y^{\mathcal{K}} \\ P(n) \leq y}} \frac{\chi(n)}{n^{\alpha / 2 + it}} - \sum_{k = 0}^{\mathcal{K}} \frac{P (\alpha / 2 + it, \chi, y)^{k}}{k!} \biggr|,
\]
where $\mathcal{K} = \lfloor \frac{\log \min \{ x^{3/2}, q \}}{\log y} \rfloor$. To handle the first term, we apply the triangle inequality and make use of the condition that $|P (\alpha/2 + it, \chi, y)| \leq V = u/3$, delivering a bound of $\ll e^{u/3}$ by completing the sum. This is $\ll \zeta (\alpha, y)^{1/2} e^{-u/6}$ (which follows immediately from the fact that $u \leq \zeta(\alpha, y) + c$ for some absolute constant $c>0$, see~\eqref{equ: u upper bound log zeta}), and seeing as the bound we are aiming for is always $\gg \zeta (\alpha, y)^{1/2} e^{-u/32}$, this is admissible. It remains to handle the second term in the previous display. We bound this using Lemma~\ref{l: good generating function}, giving
\begin{multline*}
\E_{\chi (q)} \biggl| \sum_{\substack{n \leq y^{\mathcal{K}} \\ P(n) \leq y}} \frac{\chi(n)}{n^{\alpha / 2 + it}} - \sum_{k = 0}^{\mathcal{K}} \frac{P (\alpha / 2 + it, \chi, y)^{k}}{k!} \biggr| \\ 
\ll \zeta (\alpha, y)^{1/2} \biggl(y^{(1/2 - \alpha)/2} + \exp \biggl\{ - \frac{u}{8} \biggl( \frac{\mathcal{K}}{u} -1 \biggr)^2 \biggr\} \biggr) .
\end{multline*}
Then applying the fact that $\mathcal{K}/u = \frac{\log \min\{ q, x^{3/2} \}}{\log x} + O(1/u)$, we complete the proof of Proposition~\ref{p: L1 norm}.
\end{proof}

\end{document}
